% =====================================================================
% Figure 2 — §V-H AMRF α ablation.
%
% Drop-in include for the AAAI 2027 paper main body.  Generated by
% safe-rl-2027/experiments/diagnostics/D7_amrf_alpha_analysis/plot_d7.py
% from 3 D7 instrumented_eval.npz files (Goal/Push/MGoal × seed 0 ×
% 20 episodes ≈ 45,472 per-step samples).  Native canvas is 8.6 × 3.4
% in — two-column-figure width.
%
% Usage in the body of §V-H:
%
%   % =====================================================================
% Figure 2 — §V-H AMRF α ablation.
%
% Drop-in include for the AAAI 2027 paper main body.  Generated by
% safe-rl-2027/experiments/diagnostics/D7_amrf_alpha_analysis/plot_d7.py
% from 3 D7 instrumented_eval.npz files (Goal/Push/MGoal × seed 0 ×
% 20 episodes ≈ 45,472 per-step samples).  Native canvas is 8.6 × 3.4
% in — two-column-figure width.
%
% Usage in the body of §V-H:
%
%   % =====================================================================
% Figure 2 — §V-H AMRF α ablation.
%
% Drop-in include for the AAAI 2027 paper main body.  Generated by
% safe-rl-2027/experiments/diagnostics/D7_amrf_alpha_analysis/plot_d7.py
% from 3 D7 instrumented_eval.npz files (Goal/Push/MGoal × seed 0 ×
% 20 episodes ≈ 45,472 per-step samples).  Native canvas is 8.6 × 3.4
% in — two-column-figure width.
%
% Usage in the body of §V-H:
%
%   % =====================================================================
% Figure 2 — §V-H AMRF α ablation.
%
% Drop-in include for the AAAI 2027 paper main body.  Generated by
% safe-rl-2027/experiments/diagnostics/D7_amrf_alpha_analysis/plot_d7.py
% from 3 D7 instrumented_eval.npz files (Goal/Push/MGoal × seed 0 ×
% 20 episodes ≈ 45,472 per-step samples).  Native canvas is 8.6 × 3.4
% in — two-column-figure width.
%
% Usage in the body of §V-H:
%
%   \input{figures/fig_v_h_amrf_alpha/fig_v_h_amrf_alpha.tex}
%
% Cross-references: §V-H's two findings, the §V Prop. 4 floor, and
% the §V-A recipe's hand-set $\alpha = 0.3$.
% =====================================================================
\begin{figure*}[t]
  \centering
  \includegraphics[width=\textwidth]%
    {figures/fig_v_h_amrf_alpha/fig_v_h_amrf_alpha.pdf}
  \caption{%
    \textbf{AMRF learns $\alpha(s)$ that is uncorrelated with state;
    the recipe's fixed $\alpha=0.3$ sits in the same plateau.}
    Aggregated over 45{,}472 per-step samples from 20-episode
    instrumented eval on Goal/Push/MultiGoal.
    \emph{Left:} histogram of the learned $\alpha(s)$.  Three vertical
    references: the Prop.~4 floor $\Delta t^{2}v_{\max}/r \approx
    0.05$ (green dashed), Ours' hand-set $\alpha=0.3$ (orange dotted,
    §V-A), and AMRF's empirical mean $0.87$ (red solid).  The grey
    band marks the empirical plateau $\alpha \in [0.3, 1.5]$ that
    both methods occupy.
    \emph{Right:} $\alpha(s)$ vs distance to nearest hazard, with
    scatter (subsampled to 4{,}000 points) and binned-median curve
    $\pm$ IQR.  Spearman $\rho=-0.014$ is statistically zero
    ($|\rho|<0.05$); the binned median is flat across hazard
    distances.  Conclusion: the learned $\alpha$ varies but does
    not encode a state-aware safety signal — supporting §V-H's
    claim that Prop.~4's $\alpha$ floor, not the precise value
    above it, is the load-bearing prescription.%
  }
  \label{fig:v_h_amrf_alpha}
\end{figure*}

%
% Cross-references: §V-H's two findings, the §V Prop. 4 floor, and
% the §V-A recipe's hand-set $\alpha = 0.3$.
% =====================================================================
\begin{figure*}[t]
  \centering
  \includegraphics[width=\textwidth]%
    {figures/fig_v_h_amrf_alpha/fig_v_h_amrf_alpha.pdf}
  \caption{%
    \textbf{AMRF learns $\alpha(s)$ that is uncorrelated with state;
    the recipe's fixed $\alpha=0.3$ sits in the same plateau.}
    Aggregated over 45{,}472 per-step samples from 20-episode
    instrumented eval on Goal/Push/MultiGoal.
    \emph{Left:} histogram of the learned $\alpha(s)$.  Three vertical
    references: the Prop.~4 floor $\Delta t^{2}v_{\max}/r \approx
    0.05$ (green dashed), Ours' hand-set $\alpha=0.3$ (orange dotted,
    §V-A), and AMRF's empirical mean $0.87$ (red solid).  The grey
    band marks the empirical plateau $\alpha \in [0.3, 1.5]$ that
    both methods occupy.
    \emph{Right:} $\alpha(s)$ vs distance to nearest hazard, with
    scatter (subsampled to 4{,}000 points) and binned-median curve
    $\pm$ IQR.  Spearman $\rho=-0.014$ is statistically zero
    ($|\rho|<0.05$); the binned median is flat across hazard
    distances.  Conclusion: the learned $\alpha$ varies but does
    not encode a state-aware safety signal — supporting §V-H's
    claim that Prop.~4's $\alpha$ floor, not the precise value
    above it, is the load-bearing prescription.%
  }
  \label{fig:v_h_amrf_alpha}
\end{figure*}

%
% Cross-references: §V-H's two findings, the §V Prop. 4 floor, and
% the §V-A recipe's hand-set $\alpha = 0.3$.
% =====================================================================
\begin{figure*}[t]
  \centering
  \includegraphics[width=\textwidth]%
    {figures/fig_v_h_amrf_alpha/fig_v_h_amrf_alpha.pdf}
  \caption{%
    \textbf{AMRF learns $\alpha(s)$ that is uncorrelated with state;
    the recipe's fixed $\alpha=0.3$ sits in the same plateau.}
    Aggregated over 45{,}472 per-step samples from 20-episode
    instrumented eval on Goal/Push/MultiGoal.
    \emph{Left:} histogram of the learned $\alpha(s)$.  Three vertical
    references: the Prop.~4 floor $\Delta t^{2}v_{\max}/r \approx
    0.05$ (green dashed), Ours' hand-set $\alpha=0.3$ (orange dotted,
    §V-A), and AMRF's empirical mean $0.87$ (red solid).  The grey
    band marks the empirical plateau $\alpha \in [0.3, 1.5]$ that
    both methods occupy.
    \emph{Right:} $\alpha(s)$ vs distance to nearest hazard, with
    scatter (subsampled to 4{,}000 points) and binned-median curve
    $\pm$ IQR.  Spearman $\rho=-0.014$ is statistically zero
    ($|\rho|<0.05$); the binned median is flat across hazard
    distances.  Conclusion: the learned $\alpha$ varies but does
    not encode a state-aware safety signal — supporting §V-H's
    claim that Prop.~4's $\alpha$ floor, not the precise value
    above it, is the load-bearing prescription.%
  }
  \label{fig:v_h_amrf_alpha}
\end{figure*}

%
% Cross-references: §V-H's two findings, the §V Prop. 4 floor, and
% the §V-A recipe's hand-set $\alpha = 0.3$.
% =====================================================================
\begin{figure*}[t]
  \centering
  \includegraphics[width=\textwidth]%
    {figures/fig_v_h_amrf_alpha/fig_v_h_amrf_alpha.pdf}
  \caption{%
    \textbf{AMRF learns $\alpha(s)$ that is uncorrelated with state;
    the recipe's fixed $\alpha=0.3$ sits in the same plateau.}
    Aggregated over 45{,}472 per-step samples from 20-episode
    instrumented eval on Goal/Push/MultiGoal.
    \emph{Left:} histogram of the learned $\alpha(s)$.  Three vertical
    references: the Prop.~4 floor $\Delta t^{2}v_{\max}/r \approx
    0.05$ (green dashed), Ours' hand-set $\alpha=0.3$ (orange dotted,
    §V-A), and AMRF's empirical mean $0.87$ (red solid).  The grey
    band marks the empirical plateau $\alpha \in [0.3, 1.5]$ that
    both methods occupy.
    \emph{Right:} $\alpha(s)$ vs distance to nearest hazard, with
    scatter (subsampled to 4{,}000 points) and binned-median curve
    $\pm$ IQR.  Spearman $\rho=-0.014$ is statistically zero
    ($|\rho|<0.05$); the binned median is flat across hazard
    distances.  Conclusion: the learned $\alpha$ varies but does
    not encode a state-aware safety signal — supporting §V-H's
    claim that Prop.~4's $\alpha$ floor, not the precise value
    above it, is the load-bearing prescription.%
  }
  \label{fig:v_h_amrf_alpha}
\end{figure*}
